% ===== 原文 pp.218–231（PDF p.233–246）=====
% 第 5 章 §5.4 末（pp.218–221）、§5.5（pp.222–231）
% tag 区间：5.18–5.33（(5.18) p.219；(5.19)(5.20) p.220；(5.21) p.222；(5.22) p.225；
%   (5.23)–(5.25) p.226；(5.26)–(5.31) p.227；(5.32)(5.33) p.228）
% 本块不写章标题。§5.4 标题由前块给出，本块自 p.218 中段续写，不重开小节。
% 本块写出的小节/三级标题：\subsection{周期积分的计算}（p.222）、
%   \subsubsection{积分的无关性}（p.221）、\subsubsection{原始积分}（p.222）、
%   \subsubsection{需要证明什么}（p.225）、\subsubsection{渐近行为}（p.228）、
%   \subsubsection{跨墙}（p.231）。
% 图：Figure 18（p.219）、Figure 19（p.230）、Figure 20（p.230）。
% [原文 p.218]

为证明对每个水平接缝 $\theta_{\gamma_j,\mathfrak{D}}=\mathrm{Id}$，我们按如下方式进行。对 $I\subseteq\{1,\ldots,k\}$，定义
\[
  \operatorname{Ideal}(I):=\langle u_i \mid i\notin I\rangle\subseteq \C[T_\Sigma]\otimes_\C R_k[[y_0]].
\]
我们使用归纳法，证明

\para{断言。} 对 $k'\ge 0$ 与 $\#I=k'$，有
\[
  \theta_{\gamma_j,\mathfrak{D}}\equiv \mathrm{Id}\pmod{\operatorname{Ideal}(I)}
\]
对每个水平接缝 $j$ 成立。

$k'=0$ 的初始情形是平凡的，因为所有自同构在模理想 $(u_1,\ldots,u_k)$ 下都是恒等。故设断言对所有 $k''<k'$ 已成立。固定一个满足 $\#I=k'$ 的集合 $I$。在 $M_\R\times L$ 上固定一个定向，于是若给任一接缝 $j$ 指定定向，就确定了绕 $j$ 的环路 $\gamma_j$ 的定向。我们希望研究水平接缝 $j$ 对应的 $\theta_{\gamma_j,\mathfrak{D}}$。

注意，由于水平 $j$ 的 $\theta_{\gamma_j,\mathfrak{D}}$ 只涉及与平行射线相伴随的自同构的复合，我们实际上可以写成
\[
  \theta_{\gamma_j,\mathfrak{D}}(z^{m'})=f_j^{\langle n_j,r(m')\rangle}z^{m'}
\]
其中 $n_j\in N$ 是本原的，且在 $j$ 的切空间上为零。此外，
\[
  f_j\in \C[\{m\in T_\Sigma \mid r(m)\ \text{与}\ j\ \text{相切}\}]\otimes_\C R_k[[y_0]].
\]
注意 $f_j$ 依赖于 $n_j$ 的符号选取。设我们已一致地作此选取，即若两个接缝 $j$ 与 $j'$ 有相同的切空间，则 $n_j=n_{j'}$。

我们需要证明 $f_j\equiv 1\pmod{\operatorname{Ideal}(I)}$。固定某个 $m\in T_\Sigma$。对每个水平接缝 $j$，设 $f_j$ 模 $\operatorname{Ideal}(I)$ 中含 $z^m$ 的项为 $c_{m,j}z^m$。于是 $c_{m,j}=\bar c_{m,j}\prod_{i\in I}u_i$，其中 $\bar c_{m,j}\in\C$，因为由归纳假设，对任意 $I'\subsetneq I$ 都有 $f_j\equiv 1\pmod{\operatorname{Ideal}(I')}$。

我们首先观察到：若 $\bar c_{m,j}\ne 0$，则 $r(m)\ne 0$。事实上，注意 $f_j$ 是形如 $(1+c_{m'}z^{m'})^{\pm 1}$ 的因子之积，其中按 $\theta_{\gamma_j,\mathfrak{D}}$ 的构造有 $r(m')\ne 0$。于是 $\log f_j$ 是形如 $\pm\log(1+c_{m'}z^{m'})$ 的表达式之和。用（有限的）Taylor 级数将其展开后可见，$\log f_j=\sum\pm c_{m'}z^{m'}$，且此和中出现的每个 $m'$ 都满足 $r(m')\ne 0$。另一方面，模 $\operatorname{Ideal}(I)$ 有 $\log f_j=\sum_{m\in T_\Sigma}c_{m,j}z^m$。故当 $r(m)=0$ 时 $c_{m,j}=0$，从而当 $r(m)=0$ 时 $\bar c_{m,j}=0$。

因此我们将固定满足 $r(m)\ne 0$ 的 $m\in T_\Sigma$，并专注于证明对所有水平接缝 $j$ 都有 $\bar c_{m,j}=0$。这里我们同时把 $j$ 为垂直接缝的情形纳入：对垂直接缝置 $\bar c_{m,j}=0$。注意 $\bar c_{m,j}$ 依赖于 $j$ 的定向。改变 $j$ 的定向会改变 $\gamma_j$ 的方向，从而把 $f_j$ 换成 $f_j^{-1}$，这使 $\bar c_{m,j}$ 变号。因此我们可以把
\[
  j\mapsto \bar c_{m,j}
\]
视为一维单纯复形 $\operatorname{Sing}(\mathfrak{D})$ 的一个 $1$-链。这里 $j$ 的定向选取是隐含的。

\para{子断言。} $j\mapsto\bar c_{m,j}$ 是一个 $1$-循环。

% [原文 p.219]
\begin{figure}[htbp]
\centering
\includegraphics[width=0.62\textwidth]{f5_18.png}
\figtitle{18}{图中所示的环路依次绕过 $x_1,x_2,x_3,x_4$，它在 $S\setminus\{x_1,\ldots,x_4\}$ 中可缩。}
\end{figure}

我们需要在每个间隙处验证 $1$-循环条件，故设 $(Q,P)\in\operatorname{Interstices}(\mathfrak{D})$。考虑 $M_\R\times L$ 中以 $(Q,P)$ 为中心的小二维球面 $S$。再设 $x_1,\ldots,x_s\in S$ 是互不相同的点，使得
\[
  \{x_1,\ldots,x_s\}=\bigcup_{j\in\operatorname{Joints}(\mathfrak{D})}j\cap S.
\]
取基点 $y\in S$，$y\notin\operatorname{Supp}(\mathfrak{D})$。我们可以选取 $S$ 中绕 $x_1,\ldots,x_s$ 的小逆时针环路 $\gamma_1,\ldots,\gamma_s$ 以及连接 $y$ 与 $\gamma_i$ 的基点的道路 $\beta_i$，使得
\[
  \beta_1\gamma_1\beta_1^{-1}\cdots\beta_s\gamma_s\beta_s^{-1}=1
\]
在 $\pi_1(S\setminus\{x_1,\ldots,x_s\},y)$ 中成立；见图 18。由于 $\theta_{\gamma,\mathfrak{D}}$ 只依赖于 $\gamma$ 在 $(M_\R\times L)\setminus\operatorname{Sing}(\mathfrak{D})$ 中的同伦型，我们得到等式
\begin{equation}
  \theta_{\beta_s}^{-1}\circ\theta_{\gamma_s}\circ\theta_{\beta_s}\circ\cdots\circ\theta_{\beta_1}^{-1}\circ\theta_{\gamma_1}\circ\theta_{\beta_1}=\mathrm{Id}. \tag{5.18}
\end{equation}
这里我们略去了下标中的 $\mathfrak{D}$。

% [Tutor Note: 原文 p.219 的 $\theta$ 是反变的，即 $\theta_{\gamma\delta}=\theta_\delta\circ\theta_\gamma$；这正是 (5.18) 中各项按 $s,\ldots,1$ 逆序排列的原因。核对方法：把 (5.18) 与上方的 $\beta_1\gamma_1\beta_1^{-1}\cdots\beta_s\gamma_s\beta_s^{-1}=1$ 对照，逐因子取 $\theta$ 即得左端。因此 (5.18) 就是“绕各接缝的环路之积为恒等”这一 $1$-循环条件的自同构版本。]

我们现区分两种情形。

\para{情形 1。} 间隙 $(Q,P)$ 不满足 $Q\in\{P,P_2,\ldots,P_k\}$。则由命题 5.28，对每个绕含 $(Q,P)$ 的垂直接缝的环路 $\gamma_i$ 有 $\theta_{\gamma_i}=\mathrm{Id}$。另一方面，模 $\operatorname{Ideal}(I)$，对绕水平接缝 $j_i$ 的 $\gamma_i$，由归纳假设，$f_{j_i}$ 具有 $1+(\cdots)\prod_{i\in I}u_i$ 的形式。于是可验证 $\theta_{\gamma_i}$ 模 $\operatorname{Ideal}(I)$ 必与 $\mathbb{V}_{\Sigma,k}$ 中任一元素交换。这可以像例 5.24 那样容易地看出，只需用 $u_j\prod_{i\in I}u_i\equiv 0$
% [原文 p.220]
模 $\operatorname{Ideal}(I)$ 对任意 $j$ 成立这一事实。于是特别地，$\theta_{\gamma_i}$ 与 $\theta_{\beta_i}$ 交换。从而 (5.18) 变为
\[
  \prod\theta_{\gamma_i}\equiv\mathrm{Id}\pmod{\operatorname{Ideal}(I)}
\]
其中乘积取遍所有绕水平接缝的 $\gamma_i$。把此恒等式作用于单项式 $z^{m'}$，得到
\[
  \prod f_{j_i}^{\langle n_{j_i},r(m')\rangle}z^{m'}=z^{m'}\pmod{\operatorname{Ideal}(I)},
\]
将其展开并读出 $z^{m+m'}$ 的系数，即得恒等式
\begin{equation}
  \sum\langle n_{j_i},r(m')\rangle\bar c_{m,j_i}=0 \tag{5.19}
\end{equation}
对任意 $m'\in T_\Sigma$ 成立。现在，单项式 $z^m$ 只有在 $r(m)$ 事实上与 $j_i$ 相切时才可能出现在 $f_{j_i}$ 中，故满足 $\bar c_{m,j}\ne 0$ 且含 $(Q,P)$ 的水平接缝，恰是含于仿射直线 $(Q,P)+\R(r(m),0)$ 中的那些接缝。设 $s$ 为此直线中含 $(Q,P)$ 的接缝个数，则 $s=0,1$ 或 $2$。若 $s=0$，则无须证明。若 $s=1$，设 $j_i$ 是唯一这样的接缝，则由 (5.19) 得 $\bar c_{m,j_i}=0$。若 $s=2$，设 $j_{i_1},j_{i_2}$ 是这两个接缝，则 (5.19) 蕴含 $\bar c_{m,j_{i_1}}=\bar c_{m,j_{i_2}}$，只要 $j_{i_1}$ 与 $j_{i_2}$ 定向相同。这就说明 $1$-循环条件在 $(Q,P)$ 处成立。

\para{情形 2。} 间隙 $(Q,P)$ 满足 $Q\in\{P,P_2,\ldots,P_k\}$，比如 $Q=P_i$。当 $Q=P$ 时我们记 $i=1$。论证几乎相同，但现在有两条垂直接缝，设为端点为 $(Q,P)$ 的 $j_1$ 与 $j_2$，当 $i>1$ 时 $j_1,j_2\subseteq\{P_i\}\times L$，当 $i=1$ 时 $j_1,j_2\subseteq\{(P',P')\mid P'\in L\}\subseteq M_\R\times L$。不失一般性，可取基点 $y$ 在 $x_1$ 附近，并设 $\beta_1$ 是常值道路，从而 $\theta_{\beta_1}=\mathrm{Id}$。只要我们证明
\begin{equation}
  \theta_{\beta_2}^{-1}\circ\theta_{\gamma_2}\circ\theta_{\beta_2}\circ\theta_{\gamma_1}=\mathrm{Id}, \tag{5.20}
\end{equation}
论证就与情形 1 相同。

为此，考虑 $P'\in L$、$P'$ 接近 $P$ 但 $P'\ne P$ 时的散射图 $\mathfrak{D}(\Sigma,P',P_2,\ldots,P_k)$。由注 5.29，$\mathfrak{D}(\Sigma,P',P_2,\ldots,P_k)$ 中自 $P_i$（当 $i=1$ 时为 $P'$）发出的射线与 $W_{k-1}(P_i;P')-y_0$ 中的项一一对应，其中 $W_{k-1}(P_i;P')$ 表示用标记点 $P',\ldots,P_{i-1},P_{i+1},\ldots P_k$（当 $i=1$ 时为 $P_2,\ldots,P_k$）算出的 $W_{k-1}(P_i)$。特别地，给定 $W_{k-1}(P_i;P')-y_0$ 中的一项 $cz^m$，相应的射线上携带函数 $1+u_icw(m)z^m$，其中 $w(m)$ 是 $r(m)$ 的重数。注意若 $\gamma$ 是绕 $P_i$ 的简单环路，则这样一条射线对 $\theta_{\gamma,\mathfrak{D}(\Sigma,P',P_2,\ldots,P_k)}$ 的贡献为 $\exp(\pm X_{u_icz^m})$。这里符号只依赖于 $\gamma$ 的定向以及 $\bigwedge^2M$ 与 $\Z$ 的选定等同。由例 5.24，附着于自 $P_i$ 发出的各射线的所有自同构两两交换，因为 $u_i^2=0$，故
\[
  \theta_{\gamma,\mathfrak{D}(\Sigma,P',P_2,\ldots,P_k)}=\prod\exp(\pm X_{u_icz^m})=\exp\bigl(\pm X_{u_i(W_{k-1}(P_i;P')-y_0)}\bigr).
\]
这里乘积取遍 $W_{k-1}(P_i;P')-y_0$ 中出现的一切项 $cz^m$。此外，若 $P'\in\pi_2(j_1)\setminus\{P\}$ 且 $P''\in\pi_2(j_2)\setminus\{P\}$，则由 (5.17) 对 $k-1$ 个点归纳使用（当 $i\ne 1$），以及由定理 5.35（当 $i=1$），有
\[
  u_iW_{k-1}(P_i;P'')=u_i\theta_{\beta_2}\bigl(W_{k-1}(P_i;P')\bigr).
\]
于是由引理 5.21 推出
\[
  \theta_{\gamma_2}=\bigl(\theta_{\beta_2}\circ\theta_{\gamma_1}\circ\theta_{\beta_2}^{-1}\bigr)^{-1},
\]

% [原文 p.221]
右端最后的取逆，是因为 $\gamma_1$ 与 $\gamma_2$ 分别同伦于 $M_\R\times\{P'\}$ 与 $M_\R\times\{P''\}$ 中带相反定向的环路。这就证明了 (5.20)。随后便可如情形 1 那样完成。

至此子断言的证明完成。\hfill$\square$

为完成断言的证明、从而完成定理的证明，我们现在注意：由 $j\mapsto\bar c_{m,j}$ 给出的循环 $\sigma$ 事实上为零。的确，取定一个使 $\bar c_{m,j}\ne 0$ 的接缝，$\sigma$ 是循环这一事实蕴含：含 $j$ 的直线可以写成若干接缝 $j'$ 之并，这些 $j'$ 的定向与 $j$ 上的定向相容，且 $\bar c_{m,j'}=\bar c_{m,j}$。然而，此直线中必有一个接缝 $j'$ 沿方向 $r(m)$ 无界。但 $\mathfrak{D}$ 中含 $j'$ 的任何多面体都不可能涉及形如 $z^m$ 的单项式，因为携带单项式 $z^m$ 的射线只沿方向 $-r(m)$ 无界。于是 $0=\bar c_{m,j'}=\bar c_{m,j}$，如所愿。
\end{proof}

\subsubsection{积分的无关性}

% [Tutor Note: 原文 p.221 起为三级标题 “5.4.7. Independence of the integrals.”。本节将 §5.4 的全部分析收束为两个引理：引理 5.40 说明把生成元 $\theta=\exp(cz^{m_0}X_{r(m_0)})$（$c^2=0$）作用到指数上不改变周期积分；引理 5.41 再借助定理 5.35、5.39 与引理 5.40 断言 $\int_{\Xi}e^{W_k(Q)/\hbar}\Omega$ 与 $Q,P_1,\ldots,P_k$ 的选取无关。核对锚点：引理 5.40 中 $(u,\hbar)$ 的两个分量分别是 $\widehat{\cM}_{\Sigma,k}$ 的坐标与 $\C^\times$ 中的参数，被积函数里的 $\hbar$ 与陈述中的 $\hbar$ 是同一个量。]

\begin{lemma}
  设 $\theta\in\mathbb{V}_{\Sigma,k}$，$(u,\hbar)\in\widehat{\cM}_{\Sigma,k}\times\C^\times$，并设 $f$ 属于 $\C[T_\Sigma]\otimes_\C R_k[[y_0]]$ 中由 $u_1,\ldots,u_k$ 生成的理想。则对任一循环
  \[
    \Xi\in H_2\bigl(\kappa^{-1}(u),\operatorname{Re}(W_0(Q)/\hbar)\ll 0;\C\bigr),
  \]
  有
  \[
    \int_\Xi e^{(W_0(Q)+f)/\hbar}\Omega=\int_\Xi e^{\theta(W_0(Q)+f)/\hbar}\Omega.
  \]
\end{lemma}

\begin{proof}[证明]
只需对 $\theta=\exp(cz^{m_0}X_{r(m_0)})$ 证明本引理，其中 $m_0\in T_\Sigma$，$r(m_0)\ne 0$，$c^2=0$，因为这类元素生成 $\mathbb{V}_{\Sigma,k}$。注意若 $W_0(Q)+f=\sum_m c_mz^m$，则
\[
  \theta(W_0(Q)+f)=\sum_m c_m\bigl(z^m+\langle X_{r(m_0)},r(m)\rangle cz^{m_0+m}\bigr)
\]
且
\[
  e^{\theta(W_0(Q)+f)/\hbar}=e^{(W_0(Q)+f)/\hbar}\Bigl(1+\sum_m \hbar^{-1}cc_m\langle X_{r(m_0)},r(m)\rangle z^{m_0+m}\Bigr).
\]
再者，$\dd(z^m\operatorname{dlog}(z^{m_0}))=-\langle X_{r(m_0)},r(m)\rangle z^m\Omega$。于是
\begin{align*}
  &\bigl(e^{\theta(W_0(Q)+f)/\hbar}-e^{(W_0(Q)+f)/\hbar}\bigr)\Omega\\
  &\quad = e^{(W_0(Q)+f)/\hbar}\Bigl(\hbar^{-1}\sum_m cc_m\langle X_{r(m_0)},r(m)\rangle z^{m_0+m}\Bigr)\Omega\\
  &\quad = -\dd\bigl(cz^{m_0}e^{(W_0(Q)+f)/\hbar}\operatorname{dlog}(z^{m_0})\bigr).
\end{align*}
结论随即由 Stokes 定理以及 $e^{(W_0(Q)+f)/\hbar}$ 在 $\Xi$ 的无界部分上快速趋于零这一事实得到。
\end{proof}

\begin{lemma}
  对 $\Xi\in H_2\bigl((\widehat{\cX}_{\Sigma,k})_\kappa,\operatorname{Re}(W_0(Q)/\hbar)\ll 0\bigr)$，积分
  \[
    \int_\Xi e^{W_k(Q)/\hbar}\Omega
  \]
  与 $Q$ 及 $P_1,\ldots,P_k$ 的选取无关。
\end{lemma}

\begin{proof}[证明]
这可由定理 5.35、5.39 与引理 5.40 立即得到。
\end{proof}

\subsection{周期积分的计算}

% [原文 p.222]
本节的主要目标是计算积分
\begin{equation}
  \int_{\Xi_i} e^{W_k(Q)/\hbar}\Omega \tag{5.21}
\end{equation}
在 $\PP^2$ 的情形。在此过程中，我们将证明定理 5.15 与定理 5.18。我们沿用例 5.9 中的记号 $\Sigma,T_\Sigma,t_i,\rho_i$，以及 $m_i=r(t_i)$。

\subsubsection{原始积分}

正如我们将看到的，只要知道 $W_k(Q)$ 是什么，给出 (5.21) 的具体表达式并不十分困难。这仅仅是第一步，因为要把这个显式表达式转化为对我们有用的形式，还需要付出相当大的努力；本节余下的篇幅都将用于此事。

\begin{lemma}
  限制在 $x_0x_1x_2=1$ 上，有
  \begin{align}
    &\sum_{i=0}^2\alpha^i\int_{\Xi_i}e^{(x_0+x_1+x_2)/\hbar}\Omega \notag\\
    &\quad =\hbar^{-3\alpha}\Biggl(\sum_{d=0}^{\infty}\frac{\hbar^{-3d}}{(d!)^3}
      -3\alpha\sum_{d=1}^{\infty}\frac{\hbar^{-3d}}{(d!)^3}\sum_{k=1}^{d}\frac1k \notag\\
    &\quad\qquad +\frac92\alpha^2\sum_{d=1}^{\infty}\frac{\hbar^{-3d}}{(d!)^3}
      \Biggl(\Bigl(\sum_{k=1}^{d}\frac1k\Bigr)^2+\frac13\sum_{k=1}^{d}\frac1{k^2}\Biggr)\Biggr) \notag\\
    &\quad =:\hbar^{-3\alpha}\sum_{d=0}^{\infty}\hbar^{-3d}\bigl(B_0(d)+\alpha B_1(d)+\alpha^2B_2(d)\bigr),
  \end{align}
  其中最后一个等式定义了数 $B_0(d),B_1(d),B_2(d)$。
\end{lemma}

\begin{proof}[证明]
这就是命题 2.40 所给的显式表达式的展开。
\end{proof}

% [Tutor Note: 原文 p.222 引理 5.42 的外层幂次是 $\hbar^{-3\alpha}$（指数为 $-3\alpha$，非 $-3$），而括号内的级数每项为 $\hbar^{-3d}$；末行 $\hbar^{-3\alpha}\sum_d\hbar^{-3d}(\cdots)$ 与此一致，可作为核对锚点。注意 $\alpha$ 是 Chern 类参数，$B_i(d)$ 的定义即由此行给出。]

我们可以用此式按如下方式计算 (5.21)：
\[
  \int_{\Xi_i}e^{W_k(Q)/\hbar}\Omega=e^{y_0/\hbar}\int_{\Xi_i}e^{(x_0+x_1+x_2)/\hbar}e^{(W_k(Q)-W_0(Q))/\hbar}\Omega.
\]
随后可把因子 $e^{(W_k(Q)-W_0(Q))/\hbar}$ 展成 Taylor 级数；注意在任一展开项中，$W_k(Q)-W_0(Q)$ 里的每个单项式至多出现一次，因为它的系数平方为零；因此这一展开相当简单而且是有限项。于是我们只需计算（在 $x_0x_1x_2=e^{y_1}$ 之下）
\[
  \sum_{i=0}^2\alpha^i\int_{\Xi_i}e^{(x_0+x_1+x_2)/\hbar}x_0^{n_0}x_1^{n_1}x_2^{n_2}\Omega.
\]

\begin{lemma}
  设 $x_0x_1x_2=e^{y_1}$，则
  \[
    \sum_{i=0}^2\alpha^i\int_{\Xi_i}e^{(x_0+x_1+x_2)/\hbar}x_0^{n_0}x_1^{n_1}x_2^{n_2}\Omega
    =\hbar^{-3\alpha}e^{\alpha y_1}\sum_{i=0}^2\psi_i(n_0,n_1,n_2)\alpha^i,
  \]
  其中
  \[
    \psi_i(n_0,n_1,n_2)=\sum_{d=0}^{\infty}D_i(d,n_0,n_1,n_2)\hbar^{-(3d-n_0-n_1-n_2)}e^{dy_1}
  \]
  % [原文 p.223]
  而 $D_i$ 如下给出。首先，
  \[
    D_0(d,n_0,n_1,n_2)=
    \begin{cases}
      \dfrac{1}{(d-n_0)!(d-n_1)!(d-n_2)!} & \text{若 } d\ge n_0,n_1,n_2,\\[2mm]
      0 & \text{其他情形}.
    \end{cases}
  \]
  其次，若 $d\ge n_0,n_1,n_2$，则
  \[
    D_1(d,n_0,n_1,n_2)=-\frac{\sum_{k=1}^{d-n_0}\frac1k+\sum_{k=1}^{d-n_1}\frac1k+\sum_{k=1}^{d-n_2}\frac1k}
    {(d-n_0)!(d-n_1)!(d-n_2)!},
  \]
  而若 $n_0,n_1\le d<n_2$，则
  \[
    D_1(d,n_0,n_1,n_2)=\frac{(-1)^{n_2-d-1}(n_2-d-1)!}{(d-n_0)!(d-n_1)!},
  \]
  若改为 $d<n_0$ 或 $d<n_1$，有类似的表达式。若 $d$ 小于 $n_0,n_1,n_2$ 中的两个，则
  \[
    D_1(d,n_0,n_1,n_2)=0.
  \]
  第三，若 $d\ge n_0,n_1,n_2$，则
  \[
    D_2(d,n_0,n_1,n_2)=\frac{\Bigl(\sum_{l=0}^{2}\sum_{k=1}^{d-n_l}\frac1k\Bigr)^2
    +\sum_{l=0}^{2}\sum_{k=1}^{d-n_l}\frac1{k^2}}{2(d-n_0)!(d-n_1)!(d-n_2)!},
  \]
  而若 $n_0,n_1\le d<n_2$，则
  \[
    D_2(d,n_0,n_1,n_2)=\frac{(-1)^{d-n_2}(n_2-d-1)!}{(d-n_0)!(d-n_1)!}
    \Bigl(\sum_{k=1}^{d-n_0}\frac1k+\sum_{k=1}^{d-n_1}\frac1k+\sum_{k=1}^{n_2-d-1}\frac1k\Bigr),
  \]
  若改为 $d<n_0$ 或 $d<n_1$，有类似的表达式。若 $n_0\le d<n_1,n_2$，则
  \[
    D_2(d,n_0,n_1,n_2)=\frac{(-1)^{n_1+n_2}(n_1-d-1)!(n_2-d-1)!}{(d-n_0)!},
  \]
  若改为 $n_1\le d<n_0,n_2$ 或 $n_2\le d<n_0,n_1$，有类似的表达式。最后，若 $d<n_0,n_1,n_2$，则
  \[
    D_2(d,n_0,n_1,n_2)=0.
  \]
\end{lemma}

\begin{proof}[证明]
考虑积分
\[
  I_i(a_0,a_1,a_2)=\int_{\Xi_i}e^{a_0x_0+a_1x_1+a_2x_2}\Omega,
\]
其中 $a_0,a_1,a_2\in\C^\times$ 且 $x_0x_1x_2=1$。则
\[
  \frac{\partial^{n_0+n_1+n_2}}{\partial a_0^{n_0}\partial a_1^{n_1}\partial a_2^{n_2}}I_i
  =\int_{\Xi_i}e^{a_0x_0+a_1x_1+a_2x_2}x_0^{n_0}x_1^{n_1}x_2^{n_2}\Omega.
\]
在 $a_0=a_1=a_2=e^{y_1/3}/\hbar$ 处求此式，并在积分中作变量替换 $x_i\mapsto x_ie^{-y_1/3}$。注意由于 $\Omega=\frac{\dd x_1\wedge\dd x_2}{x_1x_2}$，这样的变量替换不影响 $\Omega$。于是在 $x_0x_1x_2=e^{y_1}$ 之下得到
\[
  \left.\frac{\partial^{n_0+n_1+n_2}}{\partial a_0^{n_0}\partial a_1^{n_1}\partial a_2^{n_2}}I_i\right|_{a_i=e^{y_1/3}/\hbar}
  =\int_{\Xi_i}e^{(x_0+x_1+x_2)/\hbar}e^{-(n_0+n_1+n_2)y_1/3}x_0^{n_0}x_1^{n_1}x_2^{n_2}\Omega.
\]

% [原文 p.224]
另一方面，$I_i$ 可以通过作替换
\begin{align*}
  x_0&\mapsto (a_1a_2/a_0^2)^{1/3}x_0\\
  x_1&\mapsto (a_0a_2/a_1^2)^{1/3}x_1\\
  x_2&\mapsto (a_0a_1/a_2^2)^{1/3}x_2
\end{align*}
来计算，这给出
\[
  I_i(a_0,a_1,a_2)=\int_{\Xi_i}e^{(a_0a_1a_2)^{1/3}(x_0+x_1+x_2)}\Omega.
\]
于是，在引理 5.42 的公式中代入 $\hbar=(a_0a_1a_2)^{-1/3}$，便可算出 $\sum_{i=0}^2\alpha^iI_i(a_0,a_1,a_2)$。为对所得表达式求导，注意在此替换下 $\hbar^{-(3\alpha+3d)}$ 变为 $(a_0a_1a_2)^{\alpha+d}$，并且
\begin{align}
  &\left.\frac{\partial^{n_0+n_1+n_2}}{\partial a_0^{n_0}\partial a_1^{n_1}\partial a_2^{n_2}}(a_0a_1a_2)^{\alpha+d}\right|_{a_i=e^{y_1/3}/\hbar} \notag\\
  &\quad =\hbar^{-(3\alpha+3d-n_0-n_1-n_2)}e^{(\alpha+d-(n_0+n_1+n_2)/3)y_1} \notag\\
  &\quad\qquad\cdot\prod_{k=1}^{n_0}(\alpha+d-k+1)\prod_{k=1}^{n_1}(\alpha+d-k+1)\prod_{k=1}^{n_2}(\alpha+d-k+1) \notag\\
  &\quad =\hbar^{-(3\alpha+3d-n_0-n_1-n_2)}e^{(\alpha+d-(n_0+n_1+n_2)/3)y_1} \notag\\
  &\quad\qquad\cdot\bigl(C_0(d,n_0,n_1,n_2)+\alpha C_1(d,n_0,n_1,n_2)+\alpha^2C_2(d,n_0,n_1,n_2)\bigr),
\end{align}
其中最后一个等式定义了 $C_0,C_1$ 与 $C_2$。于是可见
\[
  \psi_i(n_0,n_1,n_2)=\sum_{d=0}^{\infty}\sum_{k=0}^{i}B_k(d)C_{i-k}(d,n_0,n_1,n_2)\hbar^{-(3d-n_0-n_1-n_2)}e^{dy_1},
\]
其中 $B_i$ 如引理 5.42 所定义。进一步，算出这些 $C_i$ 后可见
\[
  C_0(d,n_0,n_1,n_2)=
  \begin{cases}
    \dfrac{(d!)^3}{(d-n_0)!(d-n_1)!(d-n_2)!} & \text{若 } d\ge n_0,n_1,n_2,\\[2mm]
    0 & \text{其他情形}.
  \end{cases}
\]
若 $d\ge n_0,n_1,n_2$，则
\[
  C_1(d,n_0,n_1,n_2)=\frac{(d!)^3}{(d-n_0)!(d-n_1)!(d-n_2)!}
  \Bigl(\sum_{k=d-n_0+1}^{d}\frac1k+\sum_{k=d-n_1+1}^{d}\frac1k+\sum_{k=d-n_2+1}^{d}\frac1k\Bigr),
\]
而若 $n_0,n_1,n_2$ 中有一个大于 $d$，则
\[
  C_1(d,n_0,n_1,n_2)=\prod_{\substack{k=d-n_0+1\\ k\ne 0}}^{d}k\ \prod_{\substack{k=d-n_1+1\\ k\ne 0}}^{d}k\ \prod_{\substack{k=d-n_2+1\\ k\ne 0}}^{d}k.
\]
其他情形
\[
  C_1(d,n_0,n_1,n_2)=0.
\]

% [原文 p.225]
若 $d\ge n_0,n_1,n_2$，则
\begin{align*}
  &C_2(d,n_0,n_1,n_2)\\
  &\quad =\frac{(d!)^3}{2(d-n_0)!(d-n_1)!(d-n_2)!}
    \Biggl(\Bigl(\sum_{k=d-n_0+1}^{d}\frac1k+\sum_{k=d-n_1+1}^{d}\frac1k+\sum_{k=d-n_2+1}^{d}\frac1k\Bigr)^2\\
  &\qquad\qquad -\Bigl(\sum_{k=d-n_0+1}^{d}\frac1{k^2}+\sum_{k=d-n_1+1}^{d}\frac1{k^2}+\sum_{k=d-n_2+1}^{d}\frac1{k^2}\Bigr)\Biggr).
\end{align*}
若 $n_1,n_2\le d<n_0$，则
\begin{align*}
  &C_2(d,n_0,n_1,n_2)\\
  &\quad =\Bigl(\prod_{\substack{k=d-n_0+1\\ k\ne 0}}^{d}k\Bigr)\frac{(d!)^2}{(d-n_1)!(d-n_2)!}
    \Bigl(\sum_{\substack{k=d-n_0+1\\ k\ne 0}}^{d}\frac1k+\sum_{k=d-n_1+1}^{d}\frac1k+\sum_{k=d-n_2+1}^{d}\frac1k\Bigr).
\end{align*}
若改为 $d<n_1$ 或 $d<n_2$，有类似的表达式。若 $n_0,n_1,n_2$ 中有两个大于 $d$，则
\[
  C_2(d,n_0,n_1,n_2)=\prod_{\substack{k=d-n_0+1\\ k\ne 0}}^{d}k\ \prod_{\substack{k=d-n_1+1\\ k\ne 0}}^{d}k\ \prod_{\substack{k=d-n_2+1\\ k\ne 0}}^{d}k.
\]
最后，若 $n_0,n_1,n_2>d$，则
\[
  C_2(d,n_0,n_1,n_2)=0.
\]
经过一番繁复的计算，便得到引理中给出的系数 $D_i$ 的形式。
\end{proof}

\begin{definition}
  对 $m\in T_\Sigma$，$m=\sum_{i=0}^2n_it_i$，其中所有 $n_i\ge 0$，定义
  \begin{align*}
    \psi_i(m)&:=\psi_i(n_0,n_1,n_2)\\
    D_i(d,m)&:=D_i(d,n_0,n_1,n_2)
  \end{align*}
  以及
  \[
    |m|:=n_0+n_1+n_2.
  \]
\end{definition}

% [Tutor Note: 原文 p.224 的求导恒等式以 $\prod_{k=1}^{n_0}(\alpha+d-k+1)$ 表示下降阶乘 $(\alpha+d)(\alpha+d-1)\cdots(\alpha+d-n_0+1)$；令 $\alpha=0$ 即得 $C_0=(d!)^3/\prod(d-n_i)!$，可作核对锚点。又当 $n_i>d$ 时，原文把乘积的指标写作 $k=d-n_i+1$ 到 $d$ 并附注 $k\ne 0$；此附注不可省，因为此时区间含 $k=0$，若按普通乘积理解整项将为零。]

\subsubsection{需要证明什么}

如果我们已有 $W_k(Q)$ 的显式表达式，则上一小节的结论使我们能写出 (5.21)。一般说来，所得结果（除去对 $\Xi_0$ 的积分，这一点我们将会看到）与定理 5.18 所预测的期望形式并无明显的相似之处。必须先做一番整理，积分才开始显得正确。

不过，我们首先要着手证明定理 5.15 与 5.18；为此须明确：用上一小节算出的那些积分来表达，究竟需要证明哪个等式。

为记录将出现在 (5.21) 中的各项，我们作如下定义。

\begin{definition}
  固定一般的 $P_1,\ldots,P_k$。对一般的 $Q$，令 $S_k$（若需强调对 $Q$ 的依赖，则记作 $S_k(Q)$）为三元组 $(c,\nu,m)$ 的有限集合，其中 $c\in R_k$ 是单项式，$\nu\ge 0$ 是整数，$m\in T_\Sigma$，使得
  \begin{equation}
    e^{(W_k(Q)-W_0(Q))/\hbar}=\sum_{(c,\nu,m)\in S_k}c\hbar^{-\nu}z^m, \tag{5.22}
  \end{equation}
  且其中每一项 $c\hbar^{-\nu}z^m$ 都具有 $\hbar^{-\nu}\prod_{i=1}^{\nu}\operatorname{Mono}(h_i)$ 的形式，$h_1,\ldots,h_\nu$ 是以 $Q$ 为边界的、互不相同的 Maslov 指标为 $2$ 的热带圆盘。
\end{definition}

% [原文 p.226]
令
\[
  L_i^d=L_i^d(Q):=\sum_{(c,\nu,m)\in S_k}c\hbar^{-(3d+\nu-|m|)}D_i(d,m).
\]

现在我们可以弄清需要证明什么。下面的引理把定理 5.15 与 5.18 归结为三个等式。

\begin{lemma}
  设 $Q$ 取一般位置，$L$ 是以 $Q$ 为顶点的热带直线。则以下三个等式
  \begin{equation}
    L_0^d=\delta_{0,d}+\sum_{\nu\ge 0}\ \sum_{\substack{I\subseteq\{1,\ldots,k\}\\ I=\{i_1,\ldots,i_{3d-\nu-2}\}\\ i_1<\cdots<i_{3d-\nu-2}}}
    \langle P_{i_1},\ldots,P_{i_{3d-\nu-2}},\psi^\nu Q\rangle_{0,d}^{\mathrm{trop}}u_I\hbar^{-(\nu+2)} \tag{5.23}
  \end{equation}
  \begin{equation}
    L_1^d=\sum_{\nu\ge 0}\ \sum_{\substack{I\subseteq\{1,\ldots,k\}\\ I=\{i_1,\ldots,i_{3d-\nu-1}\}\\ i_1<\cdots<i_{3d-\nu-1}}}
    \langle P_{i_1},\ldots,P_{i_{3d-\nu-1}},\psi^\nu L\rangle_{0,d}^{\mathrm{trop}}u_I\hbar^{-(\nu+1)} \tag{5.24}
  \end{equation}
  \begin{equation}
    L_2^d=y_2\hbar\delta_{0,d}+\sum_{\nu\ge 0}\ \sum_{\substack{I\subseteq\{1,\ldots,k\}\\ I=\{i_1,\ldots,i_{3d-\nu}\}\\ i_1<\cdots<i_{3d-\nu}}}
    \langle P_{i_1},\ldots,P_{i_{3d-\nu}},\psi^\nu M_\R\rangle_{0,d}^{\mathrm{trop}}u_I\hbar^{-\nu} \tag{5.25}
  \end{equation}
  蕴含定理 5.15 与 5.18。
\end{lemma}

\begin{proof}[证明]
让我们明确：为证明定理 5.15 与 5.18，需要证明什么。若对 $0\le i\le 2$ 记 $y_2=\sum_{i=1}^{k}u_i$，并令
\[
  K_i^{\mathrm{trop}}=\sum_{d\ge 1}\sum_{\nu\ge 0}\langle T_2^{3d+i-2-\nu},\psi^\nu T_{2-i}\rangle_{0,d}^{\mathrm{trop}}\hbar^{-(\nu+2)}e^{dy_1}
  \frac{y_2^{3d+i-2-\nu}}{(3d+i-2-\nu)!},
\]
则
\begin{align*}
  J_0^{\mathrm{trop}}&=e^{y_0/\hbar}\bigl(1+K_0^{\mathrm{trop}}\bigr),\\
  J_1^{\mathrm{trop}}&=e^{y_0/\hbar}\bigl(\hbar^{-1}y_1(1+K_0^{\mathrm{trop}})+K_1^{\mathrm{trop}}\bigr),\\
  J_2^{\mathrm{trop}}&=e^{y_0/\hbar}\Bigl(\frac{\hbar^{-2}y_1^2}{2}(1+K_0^{\mathrm{trop}})+\hbar^{-1}y_1K_1^{\mathrm{trop}}+\hbar^{-1}y_2+K_2^{\mathrm{trop}}\Bigr).
\end{align*}
我们希望把这些表达式与经由 $\Xi_0,\Xi_1,\Xi_2$ 上的周期积分所得的表达式作比较。考虑 (5.5) 中定义的 $\varphi_i$。利用引理 5.43 与 $e^{\alpha y_1}=1+y_1\alpha+y_1^2\alpha^2/2$ 展开 (5.5) 中的积分，(5.5) 的左端即为
\begin{align*}
  &\sum_{(c,\nu,m)\in S_k}ce^{y_0/\hbar}\hbar^{-(3\alpha+\nu)}e^{\alpha y_1}\sum_{i=0}^2\psi_i(m)\alpha^i\\
  &\quad =\hbar^{-3\alpha}e^{y_0/\hbar}\sum_{(c,\nu,m)\in S_k}c\hbar^{-\nu}\sum_{i=0}^2\sum_{k=0}^{i}\frac{y_1^k}{k!}\psi_{i-k}(m)\alpha^i.
\end{align*}

% [原文 p.227]
将其与 (5.5) 的右端比较，得到
\begin{align*}
  \varphi_0&=e^{y_0/\hbar}\sum_{(c,\nu,m)\in S_k}c\hbar^{-\nu}\psi_0(m),\\
  \varphi_1&=e^{y_0/\hbar}\sum_{(c,\nu,m)\in S_k}c\hbar^{-(\nu+1)}\bigl(y_1\psi_0(m)+\psi_1(m)\bigr),\\
  \varphi_2&=e^{y_0/\hbar}\sum_{(c,\nu,m)\in S_k}c\hbar^{-(\nu+2)}\Bigl(\frac{y_1^2}{2}\psi_0(m)+y_1\psi_1(m)+\psi_2(m)\Bigr).
\end{align*}
于是，为证明 $\varphi_i=J_i^{\mathrm{trop}}$，我们需要证明以下三个等式：
\begin{equation}
  \sum_{(c,\nu,m)\in S_k}c\hbar^{-\nu}\psi_0(m)=1+K_0^{\mathrm{trop}}, \tag{5.26}
\end{equation}
\begin{equation}
  \sum_{(c,\nu,m)\in S_k}c\hbar^{-\nu}\psi_1(m)=\hbar K_1^{\mathrm{trop}}, \tag{5.27}
\end{equation}
\begin{equation}
  \sum_{(c,\nu,m)\in S_k}c\hbar^{-\nu}\psi_2(m)=\hbar^2\bigl(\hbar^{-1}y_2+K_2^{\mathrm{trop}}\bigr). \tag{5.28}
\end{equation}
然后，利用引理 5.43 中 $\psi_i$ 的展开式，若比较两端 $e^{dy_1}$ 的系数，则 (5.26)、(5.27) 与 (5.28) 分别等价于
\begin{equation}
  L_0^d=\delta_{0,d}+\sum_{\nu\ge 0}\langle T_2^{3d-\nu-2},\psi^\nu T_2\rangle_{0,d}^{\mathrm{trop}}
  \frac{y_2^{3d-\nu-2}}{(3d-\nu-2)!}\hbar^{-(\nu+2)}, \tag{5.29}
\end{equation}
\begin{equation}
  L_1^d=\sum_{\nu\ge 0}\langle T_2^{3d-\nu-1},\psi^\nu T_1\rangle_{0,d}^{\mathrm{trop}}
  \frac{y_2^{3d-\nu-1}}{(3d-\nu-1)!}\hbar^{-(\nu+1)}, \tag{5.30}
\end{equation}
\begin{equation}
  L_2^d=y_2\hbar\delta_{0,d}+\sum_{\nu\ge 0}\langle T_2^{3d-\nu},\psi^\nu T_0\rangle_{0,d}^{\mathrm{trop}}
  \frac{y_2^{3d-\nu}}{(3d-\nu)!}\hbar^{-\nu}. \tag{5.31}
\end{equation}
现在设我们已证明 (5.23)、(5.24) 与 (5.25)。这些等式的左端来自周期积分，故由引理 5.41 它们与 $Q$ 及 $P_1,\ldots,P_k$ 的位置无关。于是右端也与 $Q$ 及 $P_1,\ldots,P_k$ 的位置无关。特别地，一旦证明了 (5.23)、(5.24) 与 (5.25)，我们就发现不变量 $\langle T_2^{3d+i-2-\nu},\psi^\nu T_{2-i}\rangle_{0,d}^{\mathrm{trop}}$ 是良定义的，这就证明了定理 5.15；同时也证明了 (5.29)、(5.30) 与 (5.31)，因为
\[
  \frac{y_2^n}{n!}=\sum_{\substack{I\subseteq\{1,\ldots,k\}\\ \#I=n}}u_I.
\]
此外，这也说明 $\varphi_i=J_i^{\mathrm{trop}}$。特别地，对 $0\le i\le 2$ 有 $\varphi_{i,1}=y_i$。这就给出定理 5.18。
\end{proof}

% [Tutor Note: 原文 p.227 (5.29)–(5.31) 三项中第一个配对记号的下标一律为 $T_2$（即 $\langle T_2^{3d-\nu-2},\psi^\nu T_2\rangle$、$\langle T_2^{3d-\nu-1},\psi^\nu T_1\rangle$、$\langle T_2^{3d-\nu},\psi^\nu T_0\rangle$），与上页 $K_i^{\mathrm{trop}}$ 中的 $T_2^{3d+i-2-\nu}$ 及 $T_{2-i}$ 对应（$i=0,1,2$ 给出 $T_2,T_1,T_0$），可作核对锚点。另注意 (5.29) 与 (5.31) 的常数项分别为 $\delta_{0,d}$ 与 $y_2\hbar\delta_{0,d}$，$\hbar$ 不可漏。]

我们现在来细化需要计算的表达式 $L_i^d$，把它拆成若干项，各项将给出性质互不相同的贡献。

\begin{definition}
  对每个锥 $\sigma\in\Sigma$，$\sigma$ 都是锥 $K\subseteq T_\Sigma\otimes\R$ 的某个真面 $\tilde\sigma$ 在 $r$ 下的像，这里 $K$ 由 $t_0,t_1,t_2$ 生成（即第一卦限）。对 $d\ge 0$，
  % [原文 p.228]
  用 $K_d\subseteq K$ 记立方体
  \[
    K_d=\Bigl\{\sum_{i=0}^2n_it_i\ \Big|\ 0\le n_i\le d\Bigr\},
  \]
  并对 $\sigma\in\Sigma$ 定义
  \[
    \tilde\sigma_d:=(\tilde\sigma+K_d)\setminus\bigcup_{\substack{\tau\subsetneq\sigma\\ \tau\in\Sigma}}(\tilde\tau+K_d).
  \]
  这里 $+$ 表示 Minkowski 和。
\end{definition}

\begin{example}
  关于 $\tilde\sigma_d$ 我们有如下例子。设 $m=\sum_in_it_i\in K$。
  \begin{itemize}[leftmargin=2em,itemsep=1pt]
    \item 若 $\sigma=\{0\}$，则 $m\in\tilde\sigma_d$ 当且仅当 $d\ge\max\{n_0,n_1,n_2\}$。
    \item 若 $\sigma=\rho_0$，则 $m\in\tilde\sigma_d$ 当且仅当 $n_1,n_2\le d<n_0$。
    \item 若 $\sigma=\rho_1+\rho_2$，则 $m\in\tilde\sigma_d$ 当且仅当 $n_0\le d<n_1,n_2$。
  \end{itemize}
\end{example}

\begin{definition}
  对 $\sigma\in\Sigma$，定义
  \[
    L_{i,\sigma}^d=L_{i,\sigma}^d(Q):=\sum_{\substack{(c,\nu,m)\in S_k\\ m\in\tilde\sigma_d}}
    c\hbar^{-(3d+\nu-|m|)}D_i(d,m).
  \]
\end{definition}

于是有显然的

\begin{lemma}
  $L_i^d=\sum_{\sigma\in\Sigma}L_{i,\sigma}^d$。
\end{lemma}

% [Tutor Note: 原文 p.228 定义 5.47 中 $\tilde\sigma_d$ 的取并范围是 $\tau\subsetneq\sigma$（真子面，原文排作 $\subsetneq$），漏掉“真”字会得到 $\tilde\sigma_d=\emptyset$；此点可作核对锚点。又 $K_d$ 是以 $0,dt_0,dt_1,dt_2$ 为顶点的立方体，$\tilde\sigma_d$ 是其“边界层”的差集。]

\subsubsection{$L_{i,\omega}^d(Q)$ 的渐近行为}

理解 $\omega\ne\{0\}$ 时的项 $L_{i,\omega}^d(Q)$（$\{0\}$ 的情形将是初等的）的主要技巧，是理解当 $Q$ 沿 $\omega$ 内的一条射线跑向无穷时 $L_{i,\omega}^d(Q)$ 的渐近行为。事实上，下面的引理将说明这一项最终为零。随后我们便能通过理解 $Q$ 沿各条射线跑向无穷时该项如何变化，对任意 $Q$ 算出 $L_{i,\omega}^d(Q)$。

\begin{lemma}
  设 $\omega\in\Sigma$，$v\in\omega$ 非零（于是排除 $\omega=\{0\}$）。则
  \begin{equation}
    \lim_{s\to\infty}L_{i,\omega}^d(Q+sv)=0. \tag{5.32}
  \end{equation}
\end{lemma}

\begin{proof}[证明]
我们首先注意，当 $\omega\ne\{0\}$ 时，
\begin{equation}
  \text{若 } m\in\tilde\omega_d,\ \text{则 } r(m)\in\bigcup_{\substack{\sigma\supseteq\omega\\ \sigma\in\Sigma}}\operatorname{Int}(\sigma). \tag{5.33}
\end{equation}
其次，对充分大的 $s$，$Q+sv$ 落在 $M_\R\setminus\operatorname{Supp}(\mathfrak{D})$ 的一个无界连通分支 $\cC$ 中，其中 $\mathfrak{D}=\mathfrak{D}(\Sigma,P_1,\ldots,P_k)$。取 $s$ 充分大，可以假设 $\cC$ 是 $s\to\infty$ 时最后进入的分支。为证明 (5.32)，只需证明：若 $Q+sv\in\cC$，则存在凸锥 $K'\subseteq M_\R$，满足
\[
  K'\cap\bigcup_{\substack{\sigma\supseteq\omega\\ \sigma\in\Sigma}}\operatorname{Int}(\sigma)=\emptyset
\]
且使 $W_k(Q+sv)-W_0(Q+sv)$ 只含满足 $r(m)\in K'$ 的单项式 $z^m$。于是
\[
  \exp\bigl((W_k(Q+sv)-W_0(Q+sv))/\hbar\bigr)
\]
中的所有单项式 $z^m$ 都满足 $r(m)\in K'$，从而由 (5.33) 得 $m\notin\tilde\omega_d$。这就蕴含 (5.32)。

% [原文 p.229]
于是我们研究 $W_k(Q+sv)-W_0(Q+sv)$ 中出现的单项式 $z^m$，并构造具有所需性质的锥 $K'$。我们将用到 $\cC$ 的闭包 $\overline{\cC}$ 的渐近锥 $\operatorname{Asym}(\overline{\cC})$，其定义为 Hausdorff 极限 $\lim_{\epsilon\to 0}\epsilon\overline{\cC}$。注意 $M_\R\setminus\mathfrak{D}(\Sigma,P_1)$ 的连通分支是 $P_1-\operatorname{Int}(\sigma)$，其中 $\sigma$ 取遍 $\Sigma$ 的极大锥。由于 $\operatorname{Supp}(\mathfrak{D}(\Sigma,P_1))\subseteq\operatorname{Supp}\mathfrak{D}$，可见 $\operatorname{Asym}(\overline{\cC})$ 含于某个锥 $-\sigma$ 中，其中 $\sigma\in\Sigma$ 极大且 $(-\sigma)\cap\omega\ne\{0\}$。又注意若 $\overline{\cC}$ 的无界边互相平行，则 $\operatorname{Asym}(\overline{\cC})$ 是一条射线。记 $\partial_1,\partial_2$ 为 $\overline{\cC}$ 的两条无界边。

现在，对一般的 $s$，$W_k(Q+sv)$ 中的项 $cz^m$ 对应一条断线 $\beta$，其给定数据为 $-\infty=t_0<\cdots<t_p=0$、$m_j^\beta\in T_\Sigma$（如定义 5.30 所述），且 $m=m_p^\beta$。若 $-r(m)\notin\R_{\ge 0}v$，则对充分大的 $s$（此处 $+$ 表示 Minkowski 和）有
\[
  Q+sv\notin\R_{\ge 0}(-r(m))+\bigl(\partial\overline{\cC}\setminus(\partial_1\cup\partial_2)\bigr).
\]
的确，$\partial\overline{\cC}\setminus(\partial_1\cup\partial_2)$ 是有界的，故右端的渐近锥是 $\R_{\ge 0}(-r(m))$，而由假设它不含 $v$。于是，取充分大的 $s$，我们注意到 $\beta$ 不可能经由 $\partial\overline{\cC}\setminus(\partial_1\cup\partial_2)$ 最后进入 $\cC$，因为 $\beta$ 的最后一段线段沿方向 $-r(m)$。因此对充分大的 $s$，$\beta$ 必经由 $\partial_1$ 或 $\partial_2$ 之一进入 $\overline{\cC}$。以下我们无须研究 $-r(m)\in\R_{>0}v$ 的情形，因为我们构造的锥 $K'$ 总包含 $-v$。

现在可以假设，对大的 $s$，$\beta$ 从 $M_\R\setminus\operatorname{Supp}(\mathfrak{D})$ 的另一个无界连通分支 $\cC_{n-1}$ 进入 $\cC=\cC_n$。那么，$\beta$ 经过 $\cC_{n-1}$ 时附着于它的 $m_i^\beta$ 必满足 $-r(m_i^\beta)\notin\operatorname{Asym}(\overline{\cC}_{n-1})$。的确，否则 $\beta$ 不可能碰到 $\overline{\cC}_{n-1}$ 的无界边。同样，对充分大的 $s$，可类似地看出 $\beta$ 必经由 $\overline{\cC}_{n-1}$ 的另一条无界边进入 $\cC_{n-1}$，于是我们可以归纳地继续此过程，$\beta$ 只经由一列无界分支 $\cC_0,\ldots,\cC_n$ 穿过无界边。当 $\beta$ 拐折时，它总是向外拐折，如图 19 所示。由此我们作出以下两条观察：

\begin{itemize}[leftmargin=2.2em,itemsep=2pt]
  \item[(C1)] 若 $\operatorname{Asym}(\overline{\cC})$ 中对应于 $\partial_1$ 与 $\partial_2$ 的边分别由 $v_1,v_2$ 生成（可能 $v_1=v_2$），且 $\beta$ 经穿越 $\partial_i$ 进入 $\cC$，则 $-r(m)$ 落在以 $\R v_i$ 为边界、含 $\operatorname{Asym}(\overline{\cC})$ 的半平面中；否则 $\beta$ 无法到达 $\overline{\cC}$ 的内部。
  \item[(C2)] 对任意 $j$，$1\le j\le p$，$-r(m)$ 落在以 $\R r(m_j^\beta)$ 为边界、含 $v_i$ 的半平面中，其中 $v_i$ 对应于 $\beta$ 为进入 $\cC$ 而穿越的边 $\partial_i$。这由上述关于 $\beta$ 如何拐折的行为推出。
\end{itemize}

不妨设（为便于画图）$\omega=\rho_2$ 或 $\rho_1+\rho_2$，且 $\operatorname{Asym}(\overline{\cC})\subseteq-(\rho_0+\rho_1)$。见图 20。注意如图所示，必有 $v_2\in\rho_1+\rho_2$。

我们来分析 $\beta$ 的可能情形：依 $\beta$ 的初始方向是 $-m_1$、$-m_2$ 还是 $-m_0$，共有三种情形。

\para{情形 1。} $r(m_0^\beta)=m_1$。则 $\beta$ 必经由 $\partial_2$ 进入 $\cC$。由 (C1)，$-r(m)$ 落在以 $\R v_2$ 为边界、含 $\operatorname{Asym}(\overline{\cC})$ 的半平面中；由 (C2)，$-r(m)$ 落在以 $\R m_1$ 为边界、含 $\operatorname{Asym}(\overline{\cC})$ 的半平面中。于是
\[
  -r(m)\in\bigl(-\R_{\ge 0}m_1+\R_{\ge 0}v_2\bigr).
\]

\para{情形 2。} $r(m_0^\beta)=m_2$。则或者 $\R_{\ge 0}m_2\subseteq\operatorname{Asym}(\overline{\cC})$，或者 $\operatorname{Asym}(\overline{\cC})\subseteq\rho_1+\rho_2$，因为 $v\in\operatorname{Asym}(\overline{\cC})$。在第一种情形，$\beta$ 没有拐折的机会，故对应单项式 $x_2$，而它不出现在 $W_k(Q+sv)-W_0(Q+sv)$ 中。

% [原文 p.230]
\begin{figure}[htbp]
\centering
\includegraphics[width=0.80\textwidth]{f5_19.png}
\figtitle{19}{}
\end{figure}

\begin{figure}[htbp]
\centering
\includegraphics[width=0.62\textwidth]{f5_20.png}
\figtitle{20}{}
\end{figure}

在第二种情形，$\beta$ 在它穿越某条射线 $\mathfrak{d}\in\mathfrak{D}$ 的时刻 $t_1$ 拐折，该射线上 $f_{\mathfrak{d}}=1+c_{\mathfrak{d}}z^{m_{\mathfrak{d}}}$，且 $-r(m_{\mathfrak{d}})\in\operatorname{Int}(\rho_1+\rho_2)$。此时 $r(m_1^\beta)=m_2+r(m_{\mathfrak{d}})$，故推出 $-r(m_1^\beta)\in\rho_1+\rho_2$。（这里我们用了 $m_{\mathfrak{d}}$ 的整性与 $m_2=(0,1)$。）于是由 (C1) 与 (C2)，
\[
  -r(m)\in\bigl(\R_{\ge 0}m_1+\R_{\ge 0}v_1\bigr).
\]

% [原文 p.231]
\para{情形 3。} $r(m_0^\beta)=m_0$。此时 $\beta$ 必经由边 $\partial_1$ 进入 $\cC$，因为 $\operatorname{Asym}(\overline{\cC})\subseteq-(\rho_0+\rho_1)$。于是由 (C1) 与 (C2) 可见
\[
  -r(m)\in\bigl(\R_{\ge 0}(-m_0)+\R_{\ge 0}v_1\bigr).
\]

我们现在看到：若 $\R_{\ge 0}m_2\subseteq\operatorname{Asym}(\overline{\cC})$（当 $v$ 与 $m_2$ 成比例时总成立，特别地当 $\omega=\rho_2$ 时成立），则在这三种情形中只有情形 1 与情形 3 可能出现，并且事实上 $r(m),-v\in\rho_0+\rho_1$。于是 $K'=\rho_0+\rho_1$ 就是所要的锥，这证明了此情形下的断言。

若 $\R_{\ge 0}m_2\not\subseteq\operatorname{Asym}(\overline{\cC})$，则 $v$ 不与 $m_2$ 成比例，且 $\omega=\rho_1+\rho_2$。此时，上述三种情形表明 $-r(m)$ 总含于上半平面。于是 $K'$ 取下半平面即为所要的锥，这证明了此情形下的断言。
\end{proof}

% [Tutor Note: 原文 p.230 的图 19 与图 20 只有图号 “Figure 19”/“Figure 20” 而无图注文字，故 \figtitle 的第二参数留空；图内文字（$\cC_i$、$\partial_i$、$\beta$、$Q+sv$、$\operatorname{Asym}(\overline{\cC})$、$\rho_0+\rho_2$、$\rho_1+\rho_2$、$v_1,v_2$、$K'$、$m_i$）保持原图。图 20 中标出的 $K'$ 即引理 5.51 证明中所构造的锥。]

\subsubsection{$L_{i,\omega}^d(Q)$ 的跨墙}

下一步是说明 $L_{i,\omega}^d(Q)$ 如何通过跨墙公式依赖于 $Q$。把这公式与引理 5.51 结合，使我们能把 $L_{i,\omega}^d(Q)$ 描述为来自各次跨墙的贡献之和。

\begin{definition}
  设 $\mathfrak{D}=\mathfrak{D}(\Sigma,P_1,\ldots,P_k)$。设 $\cC_1,\cC_2$ 是 $M_\R\setminus\operatorname{Supp}(\mathfrak{D})$ 的两个连通分支，且 $\dim\overline{\cC}_1\cap\overline{\cC}_2=1$。设 $Q_i\in\cC_i$ 是一般点，$\gamma$ 是从 $Q_1$ 到 $Q_2$ 的道路，它只在某一时刻 $t_0$ 经过 $\operatorname{Supp}(\mathfrak{D})$，且 $\gamma(t_0)\notin\operatorname{Sing}(\mathfrak{D})$。设 $\mathfrak{d}\in\mathfrak{D}$ 是满足 $\gamma(t_0)\in\mathfrak{d}$ 的射线，$n_{\mathfrak{d}}\in N$ 是与 $\mathfrak{d}$ 正交的本原向量，且满足 $\langle n_{\mathfrak{d}},\gamma'(t_0)\rangle<0$。记 $f_{\mathfrak{d}}=1+c_{\mathfrak{d}}z^{m_{\mathfrak{d}}}$，注意
  \[
    \theta_{\gamma,\mathfrak{d}}(z^m)=z^m+c_{\mathfrak{d}}\langle n_{\mathfrak{d}},r(m)\rangle z^{m+m_{\mathfrak{d}}}.
  \]
  现取一对 $\omega\subsetneq\tau$，其中 $\omega,\tau\in\Sigma$ 且 $\dim\tau=\dim\omega+1$。注意存在唯一的指标 $j\in\{0,1,2\}$ 使 $m_j\notin\omega$ 而 $m_j\in\tau$；称此指标为 $j(\omega,\tau)$。然后定义
  \[
    L_{i,\mathfrak{d},\gamma,\omega\to\tau}^d:=\sum_{(c,\nu,m)}
    \langle n_{\mathfrak{d}},m_{j(\omega,\tau)}\rangle c_{\mathfrak{d}}cD_i(d,m+m_{\mathfrak{d}}+t_{j(\omega,\tau)})
    \hbar^{-(\nu+3d-|m+m_{\mathfrak{d}}|)},
  \]
  其中和取遍所有满足 $m+m_{\mathfrak{d}}\in\tilde\omega_d$ 但 $m+m_{\mathfrak{d}}+t_{j(\omega,\tau)}\in\tilde\tau_d$ 的 $(c,\nu,m)\in S_k(Q_1)$。若 $(c,\nu,m)\in S_k(Q_1)$ 满足此条件，则我们说项 $c\hbar^{-\nu}z^m$ 对 $L_{i,\mathfrak{d},\gamma,\omega\to\tau}^d$ 有贡献。

  定义
  \[
    L_{i,\gamma,\omega\to\tau}^d:=\sum_{\mathfrak{d}}L_{i,\mathfrak{d},\gamma,\omega\to\tau}^d,
  \]
  其中和取遍所有满足 $\gamma(t_0)\in\mathfrak{d}$ 的 $\mathfrak{d}\in\mathfrak{D}$。

  对 $M_\R\setminus\operatorname{Sing}(\mathfrak{D})$ 中满足 $\gamma(0)=Q$、$\gamma(1)=Q'$ 的任一路径 $\gamma$，取 $[0,1]$ 的一个分割 $0=t_0<t_1<\cdots<t_n=1$，使得 $\gamma|_{[t_{j-1},t_j]}$ 是上述那种连接相邻连通分支中端点的路径。然后定义
  \[
    L_{i,\gamma,\omega\to\tau}^d:=\sum_{j=1}^{n}L_{i,\gamma|_{[t_{j-1},t_j]},\omega\to\tau}^d.
  \]
\hfill$\square$
\end{definition}

% [Tutor Note: 原文 p.231 定义 5.52 末尾（定义陈述之后、“Here then are...” 之前）排有一个证毕符号 $\square$，显系误置（定义无待证内容，且本页已有一个 $\square$ 结束引理 5.51 的证明）；译文照原文排出该符号（与第 3 章若干定义末尾的 $\square$ 一致）。另注意本定义中 $i$ 与 $\omega\to\tau$ 的下标顺序及 $\hbar^{-(\nu+3d-|m+m_{\mathfrak{d}}|)}$ 的指数符号，可与 (5.29)–(5.31) 对照。]

下面便是相关的跨墙公式。
